\documentclass[tikz,border=10pt]{standalone}
\usepackage{amsmath,amssymb}
\usepackage{tikz}
\usetikzlibrary{arrows,automata,backgrounds,calendar,chains,decorations,decorations.pathreplacing,decorations.text,fit,matrix,mindmap,patterns,petri,plotmarks,positioning,shadows,shapes,shapes.callouts,shapes.geometric,shapes.misc,spy,trees,calc,intersections,tikzmark,arrows.meta}
\usepackage{xcolor}
\usepackage{pgfplots}
\pgfplotsset{compat=1.16}

% Common math macros from the book
\def\x{\mathbf{x}}
\def\y{\mathbf{y}}
\def\z{\mathbf{z}}
\def\A{\mathbf{A}}
\def\b{\mathbf{b}}
\def\c{\mathbf{c}}
\def\w{\mathbf{w}}
\def\v{\mathbf{v}}
\def\u{\mathbf{u}}
\def\e{\mathbf{e}}
\def\zero{\mathbf{0}}
\def\R{\mathbb{R}}
\def\N{\mathbb{N}}
\def\Z{\mathbb{Z}}
\def\Q{\mathbb{Q}}
\def\st{\text{s.t.}}
\def\vect#1{\mathbf{#1}}
\def\xLP{x^{\mathrm{LP}}}
\def\zLP{z^{\mathrm{LP}}}
\def\xIP{x^{\mathrm{IP}}}

% Common node styles from the book
\tikzset{
    main/.style={circle, draw, fill=gray!20, minimum size=1cm, font=\normalsize},
    dot/.style={circle,fill=blue,minimum size=3pt,inner sep=0pt, outer sep=-1pt},
    vertex/.style={circle,fill=black!25,minimum size=20pt,inner sep=0pt},
    edge/.style={draw,-,thick},
    directed edge/.style={draw,->,>=latex,thick},
}

% Additional macros for 3D/coordinate systems
\newcommand{\xhat}{\hat{\x}}
\newcommand{\yhat}{\hat{\y}}
\newcommand{\zhat}{\hat{\z}}
\newcommand{\ihat}{\hat{\imath}}
\newcommand{\jhat}{\hat{\jmath}}
\newcommand{\khat}{\hat{k}}

\begin{document}
\begin{tikzpicture}
    % Define the triangle vertices
    \coordinate (A1) at (0,0);
    \coordinate (B1) at (4,0);
    \coordinate (C1) at (2,3.5);

    % Draw the triangle
    \draw (A1) -- (B1) -- (C1) -- cycle;

    % Label the vertices
    \node at (A1) [below] {A};
    \node at (B1) [below] {B};
    \node at (C1) [above] {C};

    % Label the sides
    \path (A1) -- (B1) node[midway,below] {$c$};
    \path (A1) -- (C1) node[midway,left] {$b$};
    \path (B1) -- (C1) node[midway,right] {$a$};


    % Add the inequality
    \node at (6,1.75) {$\begin{aligned}
        a + b &\geq c \\
        a + c &\geq b \\
        b + c &\geq a \\
    \end{aligned}$};
    
    % Second triangle with C approaching AB
    \coordinate (A2) at (10,0);
    \coordinate (B2) at (14,0);
    \coordinate (C2) at (12,0.2); % C is very close to AB

    % Draw the triangle
    \draw (A2) -- (B2) -- (C2) -- cycle;

    % Label the vertices
    \node at (A2) [below] {A};
    \node at (B2) [below] {B};
    \node at (C2) [above] {C'};

    % Label the sides
    \path (A2) -- (B2) node[midway,below] {$c$};
    \path (A2) -- (C2) node[midway,above left] {$b'$};
    \path (B2) -- (C2) node[midway,above right] {$a'$};

    % Add the approximation
    \node at (13,1.75) {$c \approx a' + b'$};
\end{tikzpicture}
\end{document}
